% ============================================================================
% TEMPLATES ALTERNATIFS - Introduction et Contexte
% ============================================================================
% 
% Ce fichier propose plusieurs variantes de structure pour la section
% Introduction et Contexte, selon le niveau de détail et l'approche.
%
% ============================================================================

% ============================================================================
% TEMPLATE 1 : Structure concise (pour articles courts)
% ============================================================================

\section{Introduction et Contexte}

Les inondations constituent une menace majeure pour les populations et l'économie mondiale. 
L'évaluation rapide des dégâts post-désastre à partir d'images satellites est un défi critique 
pour les opérations de secours et la planification de la reconstruction.

Ce projet développe un modèle de deep learning pour classifier automatiquement les dégâts 
d'inondation au niveau bâtiment, à partir d'images satellites bi-temporelles (pré et post-désastre). 
Les contributions principales sont\: une architecture Siamese ResNet18 synchronisée, 
une data foundation anti-leakage, et l'intégration de Grad-CAM pour l'explicabilité scientifique.


% ============================================================================
% TEMPLATE 2 : Structure standard (approche équilibrée)
% ============================================================================

\section{Introduction et Contexte du Projet}

\subsection{Motivation et enjeux}
\subsubsection{Problématique des catastrophes naturelles}
Les inondations figurent parmi les catastrophes naturelles les plus coûteuses et meurtrières. 
L'accès rapide à des évaluations précises des dégâts est crucial pour\:
\begin{itemize}
    \item Les opérations de secours et l'aide humanitaire
    \item La planification de la reconstruction
    \item L'évaluation des sinistres assurés
    \item La compréhension des impacts climatiques
\end{itemize}

\subsubsection{Limitations des approches actuelles}
\begin{itemize}
    \item Analyse manuelle\: coûteuse, lente, non scalable
    \item Méthodologies empiriques\: manque de standardisation
    \item Données incomplètes\: couverture géographique limitée
\end{itemize}

\subsection{Approche proposée}
Ce projet innove en combinant\:

\textbf{1. Computer Vision bi-temporelle}\: synchronisation d'images pré/post-désastre pour détecter 
les changements réels, pas les variations naturelles.

\textbf{2. Architecture Siamese}\: apprentissage de représentations contrastives pour comparer 
les états du bâti.

\textbf{3. Explicabilité (XAI)}\: Grad-CAM pour valider scientifiquement que le modèle 
regarde les zones pertinentes.

\textbf{4. MLOps complète}\: versionning données (DVC), traçabilité modèles (MLflow), 
reproductibilité garantie.

\subsection{Objectifs}
\begin{enumerate}
    \item Développer une pipeline de classification performante et robuste
    \item Assurer la reproductibilité et la traçabilité complètes
    \item Valider l'explicabilité scientifique du modèle
    \item Préparer le déploiement industriel
\end{enumerate}


% ============================================================================
% TEMPLATE 3 : Structure détaillée universitaire (approche ambitieuse)
% ============================================================================

\section{Introduction et Contexte du Projet}

\subsection{Contexte global}
\subsubsection{Enjeux climatiques et sanitaires}
Les catastrophes naturelles liées aux changements climatiques augmentent en fréquence et intensité. 
Les inondations représentent environ 40\% des pertes économiques annuelles liées aux catastrophes 
naturelles (Ministère de la Transition Écologique, 2024). Au-delà des impacts économiques, 
les inondations menacent directement les vies humaines et les infrastructures critiques.

\subsubsection{Évaluation post-désastre}
L'évaluation rapide et précise des dégâts est fondamentale pour\:

\begin{description}
    \item[Réponse humanitaire] : identifier les zones nécessitant d'urgence
    \item[Reconstruction] : quantifier les besoins en ressources
    \item[Assurance] : valider les réclamations sinistres
    \item[Science climatique] : analyser les patterns d'intensité et de géographie
\end{description}

Cependant, les approches actuelles souffrent de graves limitations\:

\begin{itemize}
    \item \textbf{Coût} : équipes d'experts coûteux déployées sur site
    \item \textbf{Délai} : plusieurs jours/semaines pour une couverture complète
    \item \textbf{Variabilité} : évaluations subjectives et inconsistantes
    \item \textbf{Scalabilité} : impossible de couvrir les grandes géographies
\end{itemize}

\subsubsection{Imagerie satellite comme source de données}
Les images satellites haute résolution offrent\:

\begin{itemize}
    \item \textbf{Couverture} : ensemble géographique complet en quelques heures
    \item \textbf{Disponibilité} : archives historiques et données en temps réel
    \item \textbf{Objectivité} : données brutes, interprétation reproductible
    \item \textbf{Coût-efficacité} : infrastructure existante, pas de déploiement terrain
\end{itemize}

Le défi : automatiser et standardiser l'analyse de ces images.

\subsection{État de l'art}

\subsubsection{Deep Learning pour l'imagerie satellite}
Depuis AlexNet (2012) et le succès du deep learning en vision par ordinateur, plusieurs approches 
ont été adaptées à l'imagerie satellite\:

\begin{itemize}
    \item \textbf{CNNs standards} (VGG, ResNet, DenseNet) : transfert learning efficace
    \item \textbf{UNets et segmentation} : utiles pour cartographier les dégâts spatialement
    \item \textbf{Vision Transformers} : prometteurs mais exigeants en données
\end{itemize}

\subsubsection{Apprentissage bi-temporel}
Les approches multi-temporelles se classent en\:

\begin{enumerate}
    \item \textbf{Concatenation} : fusionner les images pré/post de manière brute
    \item \textbf{Siamese Networks} : extraire des features indépendantes, puis comparer
    \item \textbf{RNNs temporelles} : modéliser l'évolution temporelle
\end{enumerate}

L'approche Siamese est choisie car elle\: (a) force l'apprentissage de représentations significatives, 
(b) gère les variations non-relatifs (illumination, saisonnalité), (c) est interprétable.

\subsubsection{Explicabilité en vision}
Grad-CAM \cite{selvaraju2016grad} et Class Activation Maps (CAM) permettent de visualiser 
les zones que le réseau utilise pour sa décision. C'est crucial pour valider que le modèle 
ne repose pas sur des artefacts spurieux (nuages, ombres, route).

\subsubsection{MLOps et reproductibilité}
Les outils modernes (DVC, MLflow, Weights\&Biases) adressent le gap entre recherche et production. 
Ce projet intègre ces outils dès le départ, pas en post-hoc.

\subsection{Motivation du projet}

\subsubsection{Scientifique}
\begin{itemize}
    \item Valider l'efficacité de la synchronisation bi-temporelle pour la détection de changements
    \item Analyser la robustesse face aux données hautement déséquilibrées
    \item Démontrer le rôle de Grad-CAM dans la validation scientifique
\end{itemize}

\subsubsection{Méthodologique}
\begin{itemize}
    \item Établir un pattern reproductible de Data Foundation
    \item Proposer des bonnes pratiques d'anti-leakage pour données temporelles
    \item Intégrer MLOps dans un projet de vision par ordinateur
\end{itemize}

\subsubsection{Pratique}
\begin{itemize}
    \item Fournir un codebase production-ready
    \item Démontrer le déploiement pré-production (FastAPI, Streamlit)
    \item Préparer des bases pour la généralisation à d'autres catastrophes
\end{itemize}

\subsection{Contributions principales}

\textbf{1. Data Foundation robuste et anti-leakage}\: 
Pipeline complète pour construire des métadonnées, valider l'absence de data leakage au niveau scène, 
et générer des splits train/val/test stratifiés.

\textbf{2. Architecture Siamese ResNet18 optimisée}\: 
Modèle synchronisé pré/post avec régularisation appropriée aux données déséquilibrées.

\textbf{3. Validation scientifique par XAI}\: 
Grad-CAM appliqué systématiquement pour auditer le modèle et garantir la confiance.

\textbf{4. Pipeline MLOps end-to-end}\: 
Versioning données (DVC), tracking d'expériences (MLflow), validation automatique, prêt pour déploiement.

\subsection{Plan du mémoire}

Le document se structure en\:

\begin{description}
    \item[Section 2 : Data Foundation] Construction et validation des métadonnées, 
        gestion de l'anti-leakage, synchronisation bi-temporelle
    \item[Section 3 : Pipeline d'apprentissage] Architecture Siamese, 
        stratégies de régularisation, résultats expérimentaux
    \item[Section 4 : Explicabilité (XAI)] Grad-CAM, validation qualitative, audit du modèle
    \item[Section 5 : MLOps et versioning] DVC, MLflow, infrastructure de reproductibilité
    \item[Section 6 : Résultats et interprétation] Matrices de confusion, leçons d'ingénierie
    \item[Section 7 : Perspectives] Déploiement, monitoring, extensibilité à autres catastrophes
\end{description}


% ============================================================================
% TEMPLATE 4 : Avec équation mathématique (pour formels)
% ============================================================================

\section{Introduction et Contexte}

\subsection{Formulation du problème}

Soit $I_{pre}^{(s)}$ et $I_{post}^{(s)}$ les images satellite pré et post-désastre d'une scène $s$, 
et soit $b \in \mathbb{B}_s$ un bâtiment dans la scène. On cherche à prédire le label de dégât $y_b \in \{0,1,2,3\}$ 
(no-damage, minor, major, destroyed) à partir du couple $(I_{pre}^{(s)}[C_b], I_{post}^{(s)}[C_b])$, 
où $C_b$ est le crop défini par le bbox du bâtiment.

Notre approche maximise la probabilité a posteriori\:

$$\hat{y}_b = \arg\max_y P(y \mid I_{pre}^{(s)}[C_b], I_{post}^{(s)}[C_b], \theta)$$

où $\theta$ représentent les poids du modèle, optimisés via une fonction de perte pondérée 
pour compenser le déséquilibre des classes.

% ============================================================================
% FIN TEMPLATES
% ============================================================================
